\chapter*{Foundations, objectives and structural results}
\addcontentsline{toc}{chapter}{Foundations, objectives and structural results}
\markboth{Foundations and structural results}{Foundations and structural results}
\label{hmt:sistematica:presentacion}

Triadic Modular Holography studies the production of numerical, geometric and operator structures from oriented discrete arithmetic. Its initial object is a grid constructed from the integers one through nine. Addition and multiplication determine two evaluations of every position; modular reductions preserve remainders and quotients; a ternary unit selects operational regimes, and a transport system orders their transitions and memory. Compatible iteration of these operations constitutes the enriched state and its coinductive extension.

The governing thesis attributes to number a structure preceding its scalar evaluation. Position, orientation, carry, incidences and the chronology of operations remain in the state producing the value. The Archimedean realization is obtained by contracting compatible cylinders; the conjugate realization preserves boundary incidence. Both are readings of the same construction. In this precise sense, the real line appears as the result of double projection.

This work systematically develops that genealogy. Its main results are organized in four levels: coinductive regional generation of constants; articulation between evaluation and exceptional incidence; conservation and information-recovery laws; and Dimensional Mechanics realizations. Cardinal closure of the generated continuum and terminal realizations rest on this architecture, with the domains, operators and hypotheses established in their theorems.

\section*{Operational definition of the formal kernel}

\emph{All Possible Paths}, abbreviated APP, denotes the arithmetic support of trajectories on nine rows and nine columns. Its principal matrix is constructed by multiplying row and column indices and applying the positive digital root. If \(D_9=\{1,\ldots,9\}\), its positions are \((i,j)\in D_9^2\), and the multiplicative entry is
\[
 \Pi(i,j)=1+((ij-1)\bmod9).
\]
The additive sheet evaluates \(i+j\) on the same positions. On each sheet, the pair consisting of positive remainder and quotient restores the integer: \(n=\rho_9(n)+9q_9(n)\). The symbol \(9\) represents the zero class within the positive section. This convention makes the modular boundary visible and preserves its arithmetic lift \eqref{eq:app-restitution}.

The TRIT is the oriented ternary unit, with values \(-1,0,+1\). It acts as a reader of operational classes and a regime selector. Its arithmetic decomposition preserves the pair \(n=r_3(n)+3c_3(n)\), where \(r_3(n)\) is the balanced representative and \(c_3(n)\) its quotient. Orientation and mode selection subsequently acquire quadratic and temporal realizations, defined on their respective carriers. Composition between moduli three and nine transports both quotients; the identity \eqref{trit:eq:compatibilidad39} determines that transport exactly.

The \emph{Topological Phase Kernel}, TPK, is the topological kernel of phase. It combines selection, displacement, emission, update, lifting and return operators. A transition receives a determined state, decides the admissible operation and preserves the register needed to continue it. The enriched state thus combines APP position, active sheet, tritic orientation, remainders, carries, phase, turns and accumulated incidences. Definitions and emission rules are developed in Sections~\ref{app:construccion-explicita}, \ref{trit:capitulo} and~\ref{tpk:capitulo}.

Geometric representations are constructed on these data. APP translations determine a Cayley graph; periodic identification of boundaries produces its toroidal realization. The Kähler structure corresponds to the differential realization with declared metric and complex structure. The central block \(\left(\begin{smallmatrix}7&2\\2&7\end{smallmatrix}\right)\) admits a determinantal normalization preserving a two-dimensional Minkowski form. Each realization preserves the object and normalization from which it arises \ref{geo:kahler}--\ref{geo:lorentz}.

\section*{Regional generation and coinductive extension}

The first governing result consists in passing from the finite domain of executions to compatible generation at arbitrary depth. Initial enumeration distinguishes \(104\,976\) configurations, \(468\) distinct emissions and \(243\) admissible ternary words within the \(729=3^6\) possible words. The corresponding orbital action produces \(43\) orbits. Each cardinal counts a defined object: configurations, emissions, words or orbits.

The lifts extend the words according to the sequence
\[
 w_6\longrightarrow w_{12}\longrightarrow w_{18}
 \longrightarrow w_{24}\longrightarrow w_{30}
 \longrightarrow R_{36}\longrightarrow G_9.
\]
The first two lifts use \(L_0\); subsequent ones use \(L_1\) and preserve the register when changing regime. The boundary \(R_{36}\) opens coinductive continuation, and \(G_9\) combines the nine local realizations. The holonomy \(\Gamma_9\) acts on the enriched state. Its phase reduction is a subsequent reading.

The calendar simultaneously preserves nonadic phase and the dodecaphase register through
\[
 0\longrightarrow C_{12}\longrightarrow C_{108}
 \longrightarrow C_9\longrightarrow0.
\]
The carry of this extension determines the difference between a phase return and a state update. After nine transitions, the initial phase reappears and turn memory increases. Theorem~\ref{thm:cal-extension} and the development of holonomy~\ref{chap:sucesor-102-08} specify the composition law.

The regions of \(\pi\), \(\varphi\) and \(e\) are obtained through emitters, lifts, filters and survival conditions of the same TPK. Their compatible cylinders determine expansions to any requested depth. The five regions of \(\pi\), the \(\varphi\) axis, the joint continuations and the mirror involution are distinguished by their operations and registers. Machin, Perron--Fibonacci and exponential-semigroup characterizations enter as subsequent identifications or certificates of the constructed outputs. Generation order remains fixed by the regional operators.

\section*{Coupling constant and fine structure}

The holographic constant of arithmetic-geometric coupling is the generated dodecaphase register \(K\), with its window order, phase origin, orientation and incidences. Its determination is a central result: the conversion clock, nonadic return and phase--charge reader produce the regional descriptor; incidence and heterotypy conditions select the register in the declared terminal domain. Chapter~\ref{chap:despliegue-selector-k} separately preserves the descriptor, block repertoire, selection conditions and inversion of the signed register.

The mathematical distinctiveness of \(K\) is expressed through the operations it articulates. Its positional reading enters arithmetic closure; its incidences determine directions and projectors; its full orbit supplies a reference for operator recovery. The rational reading \(\kappa\), orbital frame \(O_K\) and memory operators receive their own definitions. This distinction reveals the effective composition between number, geometry and restitution.

The regional coordinates and dodecaphase register jointly determine the fine-structure constant. The integer route normalizes blocks with their carry and preserves the boundary when extending execution. The analytic route develops a complete function, with a recurrence for all its coefficients, and proves compatibility with the generated pre-coordinate. The analytic extension is determined within its declared class from the same state; both routes express a common coinductive fourfold relation \((\pi,\varphi,e,\alpha)\). The internal order distinguishes prior regional readings, selection of \(K\), dodecaphase closure and correlated analytic representation \ref{x_reciprocidad:subsec:alpha-lector-completo-rev08}.

\section*{Double projection and exceptional incidence}

The second governing thesis is the joint construction of the continuum. Compatible extensions constitute an inverse system; their cylinders determine Archimedean evaluation. The other projection preserves the incidences, orientations and boundary marks supporting exceptional realizations. Contraction of the reading interval and persistence of that incidence describe simultaneous properties of the same state.

The spectral construction, solenoidal balance, gauge incidence, cohomological coherence and determinant line are obtained correlatively from the common system. Their operators commute with refinements on the proved domains. This unity supplies the content of Chapters~\ref{chap:sucesor-102-05} and~\ref{chap:sucesor-102-06} and governs subsequent transports.

The exceptional chain develops Paley matrices, Witt incidences, Golay codes and lattice gluings. The marked Leech neighbour is characterized through integrality, evenness, unimodularity and root exclusion. Its extension through the lattice algebra, oscillators, cocycle and twisted sector leads to the Moonshine module. The Monster then acts on the resulting vertex operator algebra. These results are specified in Theorems~\ref{x_reciprocidad:exc:teorema-leech} and~\ref{x_reciprocidad:exc:moonshine}. Orthogonal reduction \(12\to11\to10\), isotropic reduction \(26\to24\) and momentum-winding duality are coordinated through their quotient maps. The relation with M-theory preserves this algebraic realization and the fields, actions and additional conditions of its physical realization.

\section*{Evolutionary conservation and information recovery}

Evolutionary conservation of unity is proved through maps preserving measure or norm and inverses restoring transported data. The holographic extreme and mean ratio is developed through the completion \(1000=729+271\), transfer between scales and conservation of normalized unity. Carry, transport memory, boundary incidence and the operator archive have complementary functions within that articulation.

The bilateral law establishes an identity valid for two isometric transports between Hilbert spaces, in finite or infinite dimension. In its nonadic specialization, the two channels \(Q_-\) and \(Q_+\) satisfy
\[
 9\|Q_-v\|^2+8\|Q_+v\|^2
 =17\cdot73\,\|v\|^2.
\]
The normalized channels constitute an isometry; its adjoint recovers the state and has norm one. The law also identifies which readers remain invariant and how the others transform. Recovery at arbitrary depth includes a uniform bound for the terminal component and a convergent inverse series \ref{x_reciprocidad:lib04:teo-balance}, \ref{x_reciprocidad:lib04:teo-archivo}.

Register reciprocity extends this conservation to unit-covolume deformations, dual lattices, Plücker minors and spectral forms. The flow direction arises from the generated register; regional ratios subsequently act as marks. The proofs determine both the preserved invariants and the orientations needed to invert the geometric reading.

\section*{Dimensional Mechanics realizations}

Dimensional Mechanics develops physical realizations in the order of their antecedents. The fine-structure coordinate and regional readings enter the action scale. The local determinant and its order-sixteen quotient fix the valuation producing the decade \(-34\); forward and return sections distinguish the two Planck readings. Angle and counter-angle are constructed with those coordinates and preserve their conversions between radians and degrees.

Hexadic--octadic incidence produces the channels \(90/120\). Their evaluation on the angular pair determines the Barbero--Immirzi functional; angular--Witt transport preserves the relation with mixing coordinates. The vacuum constitutive response reconstructs its electric and magnetic channels, while the gravitational section relates action to areal unity. Boltzmann transduction expresses energy per decision relative to the declared thermal section. Each constant preserves its structural definition, units and evaluation operation.

The electronic structure lies at the unique fixed centre of the atlas, \((5,5)\). The remaining material realizations are organized through routes, orbits and multisections. The atlas of eighty-one cells, fifty-six route types and thirteen multisections precede the mass readers. The composition \(\text{extension}\to\text{route}\to\text{record}\to
\text{signature}\to\text{character}\to\text{towers}\) preserves the provenance of each observable. Metrological comparisons subsequently evaluate the corresponding representations and normalizations.

\section*{Cardinal closure and regime of proofs}

Generation of the continuum receives an explicit cardinal closure. In the genealogical closure \(\mathfrak G\) defined by the treatise's rules of terms and definability, condensation places every internal rational cut at a countable stage. The code \(j(r)=\omega\beta(r)+n(r)\) distinguishes their values; encoding countable well-orders supplies the reverse injection. Theorem~\ref{ix:thm:decision-continuo-hmt} concludes
\[
 \mathfrak G_{\rm HMT}(h,m_\infty)
 \models 2^{\aleph_0}=\aleph_1,
\]
and quantifies all internal subsets of the generated line. The structure and parameters of this closure arise from the described genealogy; its representation by \(L[h,m_\infty]\) subsequently identifies the theorem's domain.

The exposition distinguishes arithmetic identities, exhaustive enumerations of finite domains, compatibility theorems and passages to the limit. Verified formalizations and numerical certificates accompany their statements with precise quantifiers and dependencies. Physical realizations add the dimensional charts, couplings and observation conditions they use. This differentiation makes each result's own strength apparent and allows its composition to be followed within a single construction.
